% ============================================================
%  part1/ch05-detector.tex  —  Chapter 5: The Detector
% ============================================================

\chapter{The Detector}
\label{ch:detector}

\epigraph{%
  Every instrument is a theory made solid.%
}{Gaston Bachelard}

\begin{WhyMatters}
  The detector chapter marks the transition from human examples
  (maps, photographs, memory) to physical measurement.
  The key move is recognising that a detector is not a
  transparent window onto a system.
  It is an interface: a two-stage projection through which
  physical states become numerical outputs.
  The kernel of this projection is not a limitation to be
  overcome but the primary object of study.
\end{WhyMatters}

\section{Measurement as two-stage projection}

Let $X$ denote a space of underlying physical states.
A detector does not usually map $X$ directly into an observed
value.
Instead, it first interacts with the system, producing an
intermediate signal space $S$, and then converts that signal
into a finite readout space $D$.
Thus measurement has the form
\[
  X \xrightarrow{F} S \xrightarrow{\pi} D.
\]
The actual observable is not $x \in X$, but
\[
  d = (\pi \circ F)(x).
\]

\begin{TechNote}[Detector signal vs.\ readout]
  The two-stage structure $X \xrightarrow{F} S \xrightarrow{\pi}
  D$ is not merely formal.
  $F$ represents physical interaction: a photon excites an
  electron, a particle ionises a gas, a gravitational wave
  stretches spacetime.
  $\pi$ represents transduction: the analogue signal is
  digitised, thresholded, or otherwise converted to a readout.
  Each stage discards different information.
  $F$ loses information about undetected degrees of freedom.
  $\pi$ loses information below the noise floor or outside the
  bandwidth.
\end{TechNote}

\begin{definition}[Detector equivalence]
  Two physical states $x_1, x_2 \in X$ are
  \emph{detector-equivalent} when
  \[
    (\pi \circ F)(x_1) = (\pi \circ F)(x_2).
  \]
  We write $x_1 \sim_D x_2$.
\end{definition}

\begin{proposition}
  Detector equivalence is an equivalence relation on $X$.
\end{proposition}

\begin{proof}
  Reflexivity: $(\pi\circ F)(x) = (\pi\circ F)(x)$.
  Symmetry: equality is symmetric.
  Transitivity: if $(\pi\circ F)(x_1) = (\pi\circ F)(x_2)$
  and $(\pi\circ F)(x_2) = (\pi\circ F)(x_3)$, then
  $(\pi\circ F)(x_1) = (\pi\circ F)(x_3)$.
\end{proof}

The detector therefore does not observe $X$.
It observes the quotient
\[
  X / {\sim_D}.
\]

\begin{ForwardRef}
  The detector equivalence class $[x]_D = (\pi\circ F)^{-1}(d)$
  is the physical version of the admissibility class introduced
  in Chapter~\ref{ch:photograph}.
  In Chapter~\ref{ch:sheaf} this becomes the fibre of the
  admissibility sheaf over an observational domain.
\end{ForwardRef}

\section{Detector bandwidth and admissible signals}

A detector is physically limited.
It has a bandwidth, a sensitivity threshold, a spatial
aperture, a response time, and a noise floor.
These limitations define an admissible signal region
\[
  S_{\mathrm{adm}} \subseteq S.
\]
Only signals satisfying $F(x) \in S_{\mathrm{adm}}$ can be
registered.

\begin{definition}[Detector-admissible state]
  A state $x \in X$ is \emph{detector-admissible} when
  $F(x) \in S_{\mathrm{adm}}$.
  The detector-admissible subset of $X$ is
  \[
    X_{\mathrm{adm}} = F^{-1}(S_{\mathrm{adm}}).
  \]
\end{definition}

\begin{proposition}
  If $S_{\mathrm{adm}} \subsetneq S$, then the detector
  ignores all states in $X \setminus X_{\mathrm{adm}}$.
\end{proposition}

\begin{proof}
  $x \notin X_{\mathrm{adm}}$ iff $F(x) \notin S_{\mathrm{adm}}$,
  i.e., the interaction signal lies outside the physically
  readable region.
  Such a state produces no detector output.
\end{proof}

The detector therefore performs two kinds of forgetting.
First, it excludes states outside its admissible domain.
Second, it identifies distinct admissible states that yield the
same output.

\principle{
  Measurement is not access.\\[4pt]
  Measurement is admissible projection.
}

\section{Reconstruction from detector output}

Given a detector output $d \in D$, the compatible physical
states are
\[
  (\pi \circ F)^{-1}(d)
  = \{x \in X : (\pi \circ F)(x) = d\}.
\]
This set is usually not a singleton.

\begin{definition}[Detector reconstruction problem]
  Given $d \in D$, find an admissible estimate
  $\hat{x} = G(d)$ such that $(\pi\circ F)(\hat{x}) = d$
  whenever possible.
  The \emph{detector reconstruction error} is
  \[
    E_D(x) = d_X\!\left(x,\, G((\pi\circ F)(x))\right).
  \]
\end{definition}

\begin{proposition}
  Exact detector reconstruction for all $x \in X$ requires
  $\pi \circ F$ to be injective.
\end{proposition}

\begin{proof}
  If $G((\pi\circ F)(x)) = x$ for all $x$, and
  $(\pi\circ F)(x_1) = (\pi\circ F)(x_2)$, then applying $G$
  gives $x_1 = x_2$.
  Contrapositive: non-injectivity implies inexact
  reconstruction.
\end{proof}

\begin{example}[Axion detector as admissibility projection]
  An axion detector searches for a dark-sector particle field
  $\Phi_{\mathrm{axion}}$ through its interaction with an
  electromagnetic field $F$.
  The interaction produces a signal $s \in S$ (a photon, a
  phonon, or a mechanical oscillation) which is then amplified
  and digitised.
  The detector cannot observe $\Phi_{\mathrm{axion}}$ directly.
  It observes $(\pi \circ F)(\Phi_{\mathrm{axion}})$.
  Observable Gaussianity of the detector output is a property
  of the projection $\pi \circ F$, not of the underlying
  axion field.
  The axion field may be non-Gaussian while the detector output
  appears Gaussian, because $\pi \circ F$ averages over many
  field degrees of freedom.
\end{example}

\begin{AdmPreview}
  The distinction between the axion field and its detector
  output is the first physical example of the principle:
  \[
    \text{observable simplicity}
    \neq
    \text{ontological simplicity.}
  \]
  This will be formalised in Chapter~\ref{ch:sheaf} using the
  admissibility sheaf, and will reappear in the cosmological
  redshift discussion in Chapter~\ref{ch:redshift}.
\end{AdmPreview}

\section{The hidden kernel}

The hidden degrees of freedom of the detector are
\[
  \ker(\pi \circ F).
\]
Two field states are observationally equivalent when
\[
  (\pi \circ F)(x_1) = (\pi \circ F)(x_2).
\]

\begin{ReconExercise}[Detector non-injectivity]
  Construct two physically distinct states $x_1 \neq x_2$
  that produce the same detector output $d$.
  \begin{enumerate}[leftmargin=2em]
    \item In a spectrometer: which aspects of a light source
      are preserved in a spectrum, and which are lost?
    \item In a thermometer: which aspects of a gas are
      preserved in a temperature reading, and which are not?
    \item In an MRI scanner: which aspects of tissue are
      preserved in the image, and which are not?
  \end{enumerate}
  For each case, identify the kernel of the measurement map
  and describe what reconstruction would require.
\end{ReconExercise}

\begin{MathEcho}
  This chapter introduced:
  \[
    X \xrightarrow{F} S \xrightarrow{\pi} D,
    \qquad
    x_1 \sim_D x_2,
    \qquad
    X_{\mathrm{adm}} = F^{-1}(S_{\mathrm{adm}}),
    \qquad
    E_D(x).
  \]
  The detector equivalence relation $\sim_D$ is the physical
  version of the projection equivalence from
  Chapter~\ref{ch:map}.
  It will appear as the fibre relation of the admissibility
  sheaf in Chapter~\ref{ch:sheaf}.
\end{MathEcho}

\WhatSurvived{Detector $X \xrightarrow{F} S \xrightarrow{\pi} D$}{%
  Detector-admissible signal\\
  Equivalence class $[x]_D$\\
  Readout $d \in D$%
}{%
  States outside $X_{\mathrm{adm}}$\\
  Sub-noise-floor variations\\
  Out-of-bandwidth modes\\
  Non-commuting observables%
}{%
  Partial; requires knowledge of $F$\\
  and additional physical constraints.%
}{%
  $\ker(\pi\circ F)$: all pairs\\
  $(x_1,x_2)$ with identical\\
  detector output.%
}
